\section{Introduction}\label{sec:intro}

Households exposed to large pork-price shocks systematically
overpredict future inflation.
Using province-level meat CPI variation from the National Bureau
of Statistics matched to four waves of the China Family Panel
Studies (2016--2022), I estimate that a one-unit increase in
the cumulative meat log-price shock generates a forecast error
of $-6.07$ percentage points under the preferred specification
with region-by-wave fixed effects (wild cluster bootstrap
$p < 0.001$, 31 province clusters); a one-standard-deviation
shock (0.28 log points) implies a forecast error of $-1.70$
percentage points.
The negative sign means that households in high-shock provinces
expected more inflation than materialised.
A signal-extraction model calibrated to the actual persistence
of meat-price shocks bounds the rational forecast error at
$|\beta_3^R| \leq \omega \approx 0.05$ per unit of the shock.
The estimated coefficient exceeds this bound by a wide margin,
and no combination of persistence assumptions within the
plausible range closes the gap.
The excess is consistent with diagnostic overreaction
\citep{Bordalo2018, Bordalo2020b}.

Why does this matter for macroeconomics?
Inflation expectations anchor wage bargaining, consumption
timing, and portfolio allocation.
If a large share of the population systematically overpredicts
inflation in response to salient relative-price movements, the
expectation wedge feeds back into real activity through
precautionary saving and consumption compression.
In China, where pork accounts for roughly 60 percent of
household meat consumption and food carries an official CPI
weight near 28 percent, the scope for salience-driven belief
distortion is larger than in economies with lower food shares
or more diversified protein markets.
The 2018--2019 African Swine Fever (ASF) episode, which reduced
the national hog herd by an estimated 30--55 percent, provides
sharp cross-province variation in the size of the salience
signal.

The identification rests on a three-equation system that
separates the belief channel from the price-level channel.
The primary test is a reduced-form regression of household
forecast errors on the province-level meat-price shock
(Equation~3).
Two decomposition equations ask whether the shock raises
household expectations (Equation~1) and whether it predicts
forward headline CPI (Equation~2).
The expectation response is positive across specifications
but imprecisely estimated under the demanding region-by-wave
fixed-effect structure; the forecast-error result is the more
powerful test because it measures the joint outcome of
expectations minus realised inflation directly and does not
require the expectation channel to be individually precise.
The treatment is measured by the NBS and merged from outside
the CFPS survey, so the household's current expectation never
appears on both sides of any equation.

I make three contributions.
First, I provide a forecast-error test of household inflation
overreaction that avoids the mechanical negative correlation
in standard revision-error designs.
The identifying variation comes from an external, supply-driven
price signal, not from within-survey belief revisions.
Second, I derive a calibrated rational benchmark from a
signal-extraction model and show that the estimated forecast
error exceeds the rational bound under any persistence
assumption.
This separates behavioural overreaction from rational forecast
errors due to transitory-shock uncertainty.
Third, I show that the overreaction is commodity-specific.
Meat shocks generate forecast errors while grain shocks of
comparable CPI weight do not, consistent with salience-driven
rather than rational updating.

The paper connects to three literatures.
Diagnostic expectations \citep{Bordalo2018, Bordalo2020b}
predict that agents overweight signals made representative
by recent experience; the forecast-error pattern I document
is a direct counterpart of this prediction.
Rational inattention \citep{Sims2003, Mackowiak2009} predicts
selective attention to informative signals; the commodity
asymmetry is consistent with disproportionate attention to
high-frequency prices.
Experience-based learning
\citep{Malmendier2016, Kuchler2019} predict that personal
purchase experience shapes expectations more than official
statistics; the meat channel operates through precisely
these signals.

Three limits apply.
The CFPS is biennial, so I cannot trace within-year adjustment
dynamics.
Province-level meat CPI conflates supply and demand movements;
ASF provides a plausible supply instrument, but the treatment
is an imperfect proxy for the supply component alone.
With 31 province clusters, inference throughout relies on
the wild cluster bootstrap.

The remainder of the paper proceeds as follows.
Section~\ref{sec:background} provides institutional background.
Section~\ref{sec:framework} develops the conceptual framework.
Section~\ref{sec:data} describes the data.
Section~\ref{sec:strategy} presents the empirical strategy.
Section~\ref{sec:results} reports main results.
Section~\ref{sec:mechanism} presents commodity placebo evidence.
Section~\ref{sec:robustness} contains robustness checks.
Section~\ref{sec:discussion} discusses alternative
interpretations and policy implications.
Section~\ref{sec:conclusion} concludes.

